\section{Restriction to a vertex block}
\label{sec:restriction}

Return to a hypergraph property $f$ on $E=E_{n,r}$ and fix the bias
$p$, with $\sigma=\sqrt{p(1-p)}$. Its Fourier coefficients use the basis of
Section~\ref{sec:fourier}, with $J=E$.
The restriction method of~\cite[Lemma~3.1 and Proposition~3.2]{Graph2026} preserves a
large permutation group after conditioning. We first obtain an
averaged estimate for restricted Fourier degrees, and then relate
those degrees to the sizes of the original supports.

For a vertex set $B\subseteq[n]$, write
\[
 F_B=\{e\in E:e\cap B\ne\varnothing\}
\]
for the edge coordinates that remain free.

\begin{lemma}[Averaged restriction estimate]\label{lem:block}
Let $r+1\le m\le n$ be integers and set $k=\sigma\log m/64$.
For $B$ chosen uniformly among the $m$-subsets of $[n]$,
\begin{equation}\label{eq:block}
 \sum_{S\subseteq E}
 \Pr_B(1\le|S\cap F_B|\le k)\widehat f(S)^2
 \le m^{-1/4}\frac{rm}{n}I_p(f).
\end{equation}
\end{lemma}

\begin{proof}
Fix $B$ and an assignment $y$ to the coordinates in $E\setminus F_B$.
Denote the resulting Boolean function on $F_B$ by $f_y$.
Every permutation supported on $B$ fixes each edge outside $F_B$
individually, and consequently fixes $y$. Thus $f_y$ is invariant
under all these permutations, regardless of the assignment.

An edge with $j$ vertices in $B$ has an orbit of size $\binom mj$:
its outside part is fixed and its inside part can be any $j$-subset
of $B$. Since $1\le j\le r<m$, this size is at least $m$.
The free product measure is invariant under the same permutations,
so influences are equal within each orbit. It follows that
\[
 I_e(f_y)\le I(f_y)/m\qquad(e\in F_B).
\]
Lemma~\ref{lem:low} therefore applies to every $f_y$.

Average its conclusion over the product law of $y$. For $A\subseteq F_B$,
the Fourier expansion of $f$ gives
\[
 \widehat{f_y}(A)=
 \sum_{D\subseteq E\setminus F_B}\widehat f(A\cup D)\chi_D(y).
\]
Orthonormality in the assigned coordinates yields
\[
 \E_y\bigl[\widehat{f_y}(A)^2\bigr]
 =\sum_{D\subseteq E\setminus F_B}\widehat f(A\cup D)^2.
\]
Also $\E_y I_e(f_y)=I_{e,p}(f)$ for every $e\in F_B$, by
conditioning the disagreement event. Consequently
\[
 \sum_{S:\,1\le|S\cap F_B|\le k}\widehat f(S)^2
 \le m^{-1/4}\sum_{e\in F_B}I_{e,p}(f).
\]
Finally average over $B$. For each fixed edge $e$, a union bound on
its $r$ vertices gives $\Pr_B(e\in F_B)\le rm/n$, proving
\eqref{eq:block}.
\end{proof}

The next lemma extends the graph capture argument
of~\cite[Lemma~3.3]{Graph2026} to identify the hypergraph supports
detected by this estimate. Here a support $S$ is itself a set of distinct $r$-edges;
its vertex union contains no isolated vertices.

\begin{lemma}[Capturing a support at one vertex]\label{lem:capture}
Let $n\ge r\ge3$ and $1\le m\le n$ be integers, let $k>0$, and put
\[
 \alpha=\frac r{r-1},\qquad L=(k/r)^\alpha.
\]
Suppose
\begin{equation}\label{eq:capture-condition}
 rL\frac{m-1}{n-1}\le\frac12.
\end{equation}
For a uniform $m$-subset $B$ of $[n]$ and every
$S\subseteq E_{n,r}$ with $1\le|S|\le L$,
\begin{equation}\label{eq:capture}
 \Pr_B(1\le|S\cap F_B|\le k)\ge\frac{m}{2n}.
\end{equation}
\end{lemma}

\begin{proof}
Fix such an $S$, write $s=|S|$, and let $V=\bigcup_{e\in S}e$
and $v=|V|$. Because $s\le\binom vr\le v^r$, the average degree
of the hypergraph with edge set $S$ is at most
\[
 \frac{rs}{v}\le r s^{(r-1)/r}\le k.
\]
Choose a vertex $u\in V$ of degree $d_S(u)$ between $1$ and $k$.
This choice depends only on $S$, before $B$ is sampled.

Conditional on $u\in B$, each other vertex lies in $B$ with
probability $(m-1)/(n-1)$. The union bound and $v\le rs\le rL$
therefore give
\begin{align*}
 \Pr_B(B\cap V=\{u\})
 &\ge\frac mn\left(1-(v-1)\frac{m-1}{n-1}\right)\\
 &\ge\frac m{2n}.
\end{align*}
On this event, the edges of $S$ meeting $B$ are exactly those
containing $u$. Thus $|S\cap F_B|=d_S(u)\in[1,k]$.
\end{proof}

The captured set need not consist of a single edge: the block meets
one support vertex, and retains every support edge incident with it.
Also the vertex $u$ may vary with $S$, since~\eqref{eq:capture}
is a separate probability bound for each nonnegative summand
in~\eqref{eq:block}.

Combining Lemmas~\ref{lem:block} and~\ref{lem:capture} gives,
whenever their hypotheses hold,
\begin{equation}\label{eq:captured-mass}
 \sum_{1\le|S|\le L}\widehat f(S)^2
 \le 2r m^{-1/4}I_p(f).
\end{equation}
If $L<1$, this conclusion holds because the sum is empty.
Otherwise, restrict the nonnegative sum in~\eqref{eq:block} to
$1\le|S|\le L$ and divide by the lower bound $m/(2n)$.
This completes the transfer from restricted degree $k$ to original
degree $(k/r)^{r/(r-1)}$.
